% U4 WS4 -- Type C: long FRQ matching the exam format, rubric visible.
\documentclass{apchem-worksheet}

\apcunit{4}
\apcunittitle{Chemical Reactions}
\apcdoctitle{Worksheet 4 \textbullet\ FRQ Practice}
\apcsubtitle{Type C \textbullet\ one long FRQ, the exam's format \textbullet\ answer, then self-score}

\begin{document}
\wsheader[The unit exam's free-response question is this shape and length.
Write everything before looking at the rubric, then grade yourself the way a
reader would --- no credit for anything not written down.]

\problem[4.2]{\textbf{Long FRQ (10 points).} A student investigates the
reaction between aqueous lead(II) nitrate and aqueous potassium iodide. She
mixes \SI{50.0}{\milli\liter} of \SI{0.100}{\Molar} \ce{Pb(NO3)2} with
\SI{50.0}{\milli\liter} of \SI{0.100}{\Molar} \ce{KI} and observes a bright
yellow solid.}

\begin{parts}
  \item Write the balanced formula equation, including states.
        \answerlines[2]{\ce{Pb(NO3)2(aq) + 2 KI(aq) -> PbI2(s) +
        2 KNO3(aq)}}
  \item Write the net ionic equation and identify the spectator ions.
        \answerlines[2]{\ce{Pb^2+(aq) + 2 I^-(aq) -> PbI2(s)}; spectators
        are \ce{K+} and \ce{NO3-}.}
  \item Determine the limiting reactant, showing the comparison that
        justifies your answer. \workspace[2.8cm]
        \keyonly{\begin{apcanswer}$n(\ce{Pb^2+}) = (0.0500)(0.100) =
        \SI{5.00e-3}{\mole}$; $n(\ce{I-}) = \SI{5.00e-3}{\mole}$.
        The ratio required is 1 \ce{Pb^2+} : 2 \ce{I-}, so iodide runs out
        first --- \ce{I-} is limiting.\end{apcanswer}}
  \item Calculate the mass of \ce{PbI2} ($M = \SI{461.0}{\gram\per\mole}$)
        produced. \answerlines[2]{$n(\ce{PbI2}) = \tfrac12(5.00\times10^{-3})
        = \SI{2.50e-3}{\mole}$; $m = 2.50\times10^{-3} \times 461.0 =
        \SI{1.15}{\gram}$.}
  \item In the box, draw a particulate representation of the solution
        \emph{after} the reaction is complete. Show at least six particles,
        label your symbols, and omit water. \workspace[4cm]
        \keyonly{\begin{apcanswer}A cluster of \ce{PbI2} solid; separated
        \ce{K+} and \ce{NO3-} ions in solution; and leftover \ce{Pb^2+}
        ions (half the original amount) still dissolved, since iodide was
        limiting.\end{apcanswer}}
  \item The student claims the yellow colour alone proves a chemical
        reaction occurred. Evaluate this claim.
        \answerlines[2]{Colour alone is weak evidence --- it could arise
        from mixing coloured solutions. The stronger evidence is that an
        insoluble \emph{solid} formed from two clear solutions, which no
        physical process explains.}
\end{parts}

\begin{apcnote}[Rubric --- score yourself, 10 points]
\textbf{(a) 1 pt} balanced equation with correct states \textbullet\
\textbf{(b) 2 pts} 1 net ionic with the 2:1 coefficient, 1 both spectators
\textbullet\ \textbf{(c) 2 pts} 1 both mole values, 1 comparison that uses
the 1:2 ratio and names \ce{I-} --- comparing raw moles earns 0 \textbullet\
\textbf{(d) 2 pts} 1 applies the 2:1 ratio to get 2.50e-3 mol, 1 mass with
units \textbullet\ \textbf{(e) 2 pts} 1 precipitate drawn as a cluster with
spectators separated, 1 leftover \ce{Pb^2+} shown \textbullet\
\textbf{(f) 1 pt} identifies precipitate formation as the stronger
evidence.
\end{apcnote}

\problem[4.9]{Short redox practice. For
\ce{2 KMnO4 + 5 H2C2O4 + 3 H2SO4 -> 2 MnSO4 + 10 CO2 + K2SO4 + 8 H2O}:}
\begin{parts}
  \item Oxidation state of Mn on each side:
        \answerblank[4cm]{$+7 \to +2$}
  \item Oxidation state of C in \ce{H2C2O4} and in \ce{CO2}:
        \answerblank[4cm]{$+3 \to +4$}
  \item Oxidizing agent and reducing agent:
        \answerblank[7.4cm]{\ce{KMnO4} oxidizing; \ce{H2C2O4} reducing}
\end{parts}

\begin{spiralstrip}{Units 1--4 \textbullet\ spiral}
  \item Geometry of \ce{SO3}: \answerblank[3.4cm]{trigonal planar}
  \item Higher bp: \ce{H2O} or \ce{H2S}? \answerblank[2.4cm]{\ce{H2O}}
  \item Moles in \SI{50.0}{\milli\liter} of \SI{0.200}{\Molar}:
        \answerblank[2.8cm]{\SI{0.0100}{\mole}}
  \item PES: peak area tells you: \answerblank[5.2cm]{electrons in that
        subshell}
  \item Empirical formula of \ce{C2H6O2}: \answerblank[2.4cm]{\ce{CH3O}}
\end{spiralstrip}

\end{document}
